\clearpage
\renewcommand{\sectioncolor}{csB}
\section{Case Study 2 — What is \texorpdfstring{$e$}{e}?}
\markboth{Case Study 2 · What is $e$?}{}
\noindent\chip{csB}{BOOK pp.\ 399--419}\;\chip{slate}{PDF pp.\ 418--438}\;\chip{csB}{DELIVERS: $e$ as a \emph{concept}, and why it is 2.71828\ldots}

\vspace{6pt}
\subsection{Demolition: exponentiation is not self-multiplication}

\begin{critbox}[title={The confusion that survives into university teaching}]
$2^{2/3}$ is \textbf{not} $2$ multiplied by itself two-thirds of a time. $2^{\pi}$ is \textbf{not} $2$
multiplied by itself $3.14159\ldots$ times.
\begin{bookquote}{book p.~404 / pdf p.~423}
\textit{The exponential mapping for a base $b$ is defined as specified above, but it is not
self-multiplication.}
\end{bookquote}
Self-multiplication is a \emph{special case that happens to coincide on the integers}, and it rests
on a \emph{different} metaphor --- \meta{Arithmetic Is Object Construction} (a whole made of equal
parts). Hence the symbol ``$2^5$'' is conceptually \textbf{ambiguous}.
\end{critbox}

\begin{center}
\begin{tikzpicture}[font=\footnotesize,
  rd/.style={rounded corners=4pt,draw=#1,line width=1pt,fill=#1!8,inner sep=8pt,
             text width=6.6cm,align=center}]
  \node[fill=ink,text=white,rounded corners=4pt,inner sep=8pt,font=\large\bfseries] (sym) at (0,1.9) {$2^5$};
  \node[font=\scriptsize,text=slate,right=5pt] at (0.6,1.9) {one symbol, two ideas};
  \node[rd=teal] (l) at (-4.4,0) {\textbf{\color{teal}Object-construction reading}\\[3pt]
     $2\cdot2\cdot2\cdot2\cdot2$\\[3pt]
     {\scriptsize a \emph{whole} made of five equal \emph{parts}}\\[2pt]
     {\scriptsize grounding metaphor: \meta{Arithmetic Is Object Construction}}};
  \node[rd=csB] (r) at (4.4,0) {\textbf{\color{csB!80!black}Exponential-mapping reading}\\[3pt]
     $\mathrm{Exp}_2(5)$\\[3pt]
     {\scriptsize ``$5$ is the \emph{input}; $2$ says what $1$ maps to''}\\[2pt]
     {\scriptsize the only reading available for $2^{\pi}$, $2^{2/3}$, $2^{i}$}};
  \draw[-{Stealth[length=5pt]},line width=1pt,draw=teal]  (sym)--(l);
  \draw[-{Stealth[length=5pt]},line width=1pt,draw=csB]   (sym)--(r);
  \node[font=\scriptsize,text=crimson,align=center] at (4.4,-1.95)
     {\textbf{Only this branch survives}\\into $e^{\pi i}$};
  \draw[crimson,-{Stealth[length=4pt]},line width=0.8pt] (4.4,-1.55)--(4.4,-1.32);
\end{tikzpicture}
\end{center}
\vspace{-4pt}
\begin{center}\begin{minipage}{0.93\textwidth}\footnotesize\color{slate}
\textbf{Figure 5.}\; The ambiguity of exponential notation \pg{404}{423}. The book proposes writing
$\mathrm{Exp}_b(x)$ precisely because the ordinary notation \emph{hides} the asymmetry between
$b$ and $x$.
\end{minipage}\end{center}

\begin{ideabox}[title={Where $b$ and $x$ actually live --- and why $b^0=1$}]
In ``$b^x$'' the two symbols are not on the same footing at all:
\[
\underbrace{b}_{\text{\color{crimson}in the \textbf{range}}}{}^{\textstyle\underbrace{x}_{\text{\color{teal}in the \textbf{domain}}}}
\qquad\text{---}\qquad
b\;\text{is merely \emph{the number that }} 1 \text{\emph{ maps onto}.}
\]
So $b^0=1$ is not a convention and not ``$b$ multiplied by itself zero times'' (which is nonsense).
It says: \textbf{the exponential sends the additive identity $0$ to the multiplicative identity $1$}
--- true for every base, which is why $b^0=1$ always.
\begin{bookquote}{book p.~405 / pdf p.~424}
\textit{The way that the exponential function is written---``$b^x$''---is misleading. It is done that
way for the purpose of ease of calculation, for those cases where self-multiplication will happen to
allow you to calculate the value of the function.}
\end{bookquote}
\end{ideabox}

\subsection{What a logarithm really is}

Not ``the inverse of a power'', but a \textbf{mapping that lets you multiply by adding} --- in modern
terms a group homomorphism $(\mathbb{R}^{+},\times)\to(\mathbb{R},+)$. Napier's labour was tabulating it;
the slide rule is the physical instrument of the same idea.

\begin{center}
\begin{tikzpicture}[font=\footnotesize,scale=0.98]
% ---------- domain line (multiplicative) ----------
\draw[csB,line width=1.5pt,-{Stealth[length=5pt]}] (0.2,2.2)--(15.5,2.2);
\node[anchor=east,font=\scriptsize\bfseries,text=csB,align=right] at (0.05,2.2)
  {DOMAIN\\{\scriptsize$(\mathbb{R}^{+},\times)$}};
\foreach \x/\l in {1.3/{$\tfrac1{100}$},3.35/{$\tfrac1{10}$},5.4/{$1$},7.45/{$10$},
                   9.5/{$100$},11.55/{$1{,}000$},13.6/{$10{,}000$}}{
  \fill[csB] (\x,2.2) circle (2pt); \node[above=4pt,font=\scriptsize] at (\x,2.2) {\l};}
% ---------- range line (additive) ----------
\draw[teal,line width=1.5pt,-{Stealth[length=5pt]}] (0.2,-0.95)--(15.5,-0.95);
\node[anchor=east,font=\scriptsize\bfseries,text=teal,align=right] at (0.05,-0.95)
  {RANGE\\{\scriptsize$(\mathbb{R},+)$}};
\foreach \x/\l in {1.3/{$-2$},3.35/{$-1$},5.4/{$0$},7.45/{$1$},
                   9.5/{$2$},11.55/{$3$},13.6/{$4$}}{
  \fill[teal] (\x,-0.95) circle (2pt); \node[below=4pt,font=\scriptsize] at (\x,-0.95) {\l};}
% ---------- mapping arrows ----------
\foreach \x in {1.3,3.35,5.4,7.45,9.5,11.55,13.6}{
  \draw[purple!45,-{Stealth[length=3.5pt]},line width=0.55pt] (\x,2.0)--(\x,-0.75);}
\node[font=\scriptsize,text=purple,fill=white,inner sep=1.5pt] at (2.33,0.62) {$\log$};
\node[font=\scriptsize,text=slate,align=center,fill=white,inner sep=2pt] at (5.4,0.62)
  {\textbf{$\log$: products $\mapsto$ sums}\\[1pt]
   $1\mapsto0$,\quad $b\mapsto1$,\quad $b^{n}\mapsto n$};
% ---------- worked example : 10 x 100 = 1000 ----------
\foreach \x in {7.45,9.5}{\draw[crimson,line width=1.1pt] (\x,2.2) circle (4.5pt);
                          \draw[crimson,line width=1.1pt] (\x,-0.95) circle (4.5pt);}
\draw[crimson,line width=1.3pt] (11.55,2.2) circle (5.5pt);
\draw[crimson,line width=1.3pt] (11.55,-0.95) circle (5.5pt);
% brace under range
\draw[decorate,decoration={brace,amplitude=4pt,mirror},crimson,line width=0.8pt]
   (7.45,-1.55)--(9.5,-1.55);
\node[font=\scriptsize,text=crimson,below=4pt] at (8.47,-1.62) {\textbf{add}: $1+2=3$};
\draw[crimson,line width=1pt,-{Stealth[length=4.5pt]}]
   (9.75,-1.95) .. controls (10.9,-2.05) and (11.55,-1.9) .. (11.55,-1.42);
% up arrow: inverse log
\draw[crimson,line width=1.1pt,-{Stealth[length=5pt]}] (12.15,-0.95) -- (12.15,2.05);
\node[font=\scriptsize,text=crimson,anchor=west] at (12.28,0.62) {$\log^{-1}$};
% top annotation
\node[font=\scriptsize,text=crimson,align=center] at (9.5,3.52)
   {\textbf{multiply}:\; $10\times100=1{,}000$};
\draw[crimson,line width=0.7pt,dashed] (9.5,3.30)--(9.5,3.06);
\draw[crimson,line width=0.7pt,dashed] (7.45,2.96)--(11.55,2.96);
\draw[crimson,line width=0.7pt,dashed] (7.45,2.96)--(7.45,2.80);
\draw[crimson,line width=0.7pt,dashed] (11.55,2.96)--(11.55,2.80);
\end{tikzpicture}
\end{center}
\vspace{-3pt}
\begin{center}\begin{minipage}{0.93\textwidth}\footnotesize\color{slate}
\textbf{Figure 6.}\; The logarithm mapping, base 10 \pg{400--402}{419--421} (the book's Figures
CS\,2.1--2.2). The four constraints listed in the middle \textbf{completely and uniquely determine
every value} of the mapping --- though, as the book notes, they do not by themselves give an
algorithm for computing it.
\end{minipage}\end{center}

\noindent The exponential is simply this mapping run backwards: $0\mapsto1$, $1\mapsto b$,
$n\mapsto b^n$, \textbf{sums $\mapsto$ products}. Cognitively, Lakoff \& Núñez read it as the
arithmetization of the metaphor \meta{Multiplication Is Addition} \pg{405}{424}.

\clearpage
\subsection{The seven-step mathematicization of change}

For a foundations student this is the single most valuable page in Part VI. Before differential
calculus can be ``the mathematics of change'', \textbf{seven} conceptual moves must already have been
made --- and by the time you meet them in a calculus course they have all been internalised
invisibly.

\begin{center}
\begin{tikzpicture}[font=\footnotesize,
  st/.style={rounded corners=3pt,draw=csB,line width=0.9pt,fill=csB!8,
             text width=7.15cm,align=left,inner sep=6pt},
  num/.style={circle,fill=csB,text=white,font=\scriptsize\bfseries,inner sep=3pt},
  dn/.style={-{Stealth[length=5pt,width=4pt]},line width=1pt,draw=csB!70}]
\node[st] (s1) at (0,0) {\textbf{Time is distance.}\; The everyday metaphor: \emph{we're coming up on
  Christmas}, \emph{a long way to the end of the project}.};
\node[st] (s2) at (0,-1.62) {\textbf{Motion $\to$ two independent dimensions}, distance and
  time-as-distance.};
\node[st] (s3) at (0,-3.24) {\textbf{\meta{Logical Independence Is Geometrical Orthogonality}} ---
  this is what makes them \emph{axes}.};
\node[st] (s4) at (0,-5.10) {\textbf{Motion becomes a function $\mathbb{R}\to\mathbb{R}$.}\; The motion
  itself does not survive: it becomes a \emph{set of pairs} $(t,l)$. Continuity is supplied
  metaphorically by Weierstrass's preservation-of-closeness.};
\node[st] (s5) at (0,-7.00) {\textbf{Average speed becomes a ratio of numbers.}\; ``Average speed, of
  course, does not occur in the world.''};
\node[st] (s6) at (0,-8.86) {\textbf{\meta{Instantaneous Speed Is Average Speed over an Interval of
  Infinitely Small Size}} --- which itself requires the \textbf{BMI}.};
\node[st] (s7) at (0,-10.72) {\textbf{\meta{Change Is Motion}}, to generalise from speed of motion to
  rate of change \emph{überhaupt}.};
\foreach \i in {1,...,7}{\node[num] at ($(s\i.west)+(-0.42,0)$) {\i};}
\foreach \a/\b in {s1/s2,s2/s3,s3/s4,s4/s5,s5/s6,s6/s7}{\draw[dn] (\a)--(\b);}
% right-hand commentary
\node[anchor=north west,text width=6.9cm,align=left,font=\scriptsize,text=slate]
   at (4.35,0.62)
   {\textbf{\color{csB!80!black}\normalsize Read the stack downward.}\\[4pt]
    Each step is a \emph{substitution}: something directly perceived is replaced by something
    calculable. Motion is perceived directly, by motion detectors in the visual system. Nothing in
    the final object --- a set of ordered pairs of real numbers --- moves.\\[7pt]
    \textbf{\color{crimson}What happens if you forget they are metaphors:}};
\node[anchor=north west,inner sep=0pt] at (4.35,-2.85)
 {\begin{minipage}{6.9cm}
  \begin{bookquote}{Russell, \textit{Principles of Mathematics} 1903, p.~473; quoted book p.~408 / pdf p.~427}
  \textit{Motion consists merely in the occupation of different places at different times, subject to
  continuity.\ldots\ There is no transition from place to place, no consecutive moment or consecutive
  position, no such thing as velocity except in the sense of a real number, which is the limit of a
  certain set of quotients.}
  \end{bookquote}
  \end{minipage}};
\node[anchor=north west,text width=6.9cm,align=left,font=\scriptsize,text=slate] at (4.35,-8.15)
   {Lakoff \& Núñez call this \textbf{``pathological literalness''} --- mistaking the output of the
    metaphor chain for the thing itself. Their diagnosis:\\[4pt]
    \textit{``His view that thought was mathematical logic blinded him to the metaphorical mechanisms
    of mind needed to go from the normal human concept of change to a conceptualization of change in
    Weierstrass's mathematics.''}\\[3pt]
    \hfill\textsc{book p.~408 / pdf p.~427}};
\end{tikzpicture}
\end{center}
\vspace{-4pt}
\begin{center}\begin{minipage}{0.93\textwidth}\footnotesize\color{slate}
\textbf{Figure 7.}\; The seven steps \pg{407--408}{426--427}. Keep this beside any rigorous
treatment of the derivative: it is an inventory of everything you silently import when you say
``calculus is the mathematics of change''.
\end{minipage}\end{center}

\clearpage
\subsection{The derivation --- and the exact location of the hinge}

Now the payoff. Differentiate a general exponential and watch where the whole of Part VI turns.

\begin{center}
\begin{tcolorbox}[enhanced,width=0.97\textwidth,colback=white,colframe=csB,boxrule=1pt,arc=3pt,
   title={\bfseries Differentiating $b^{x}$ --- the load-bearing step is coloured},
   fonttitle=\bfseries,coltitle=white,colbacktitle=csB]
\begin{align*}
\frac{d}{dx}b^{x}
 &= \lim_{\Delta x\to 0}\frac{b^{\,x+\Delta x}-b^{x}}{\Delta x} &&\text{\footnotesize derivative metaphor + BMI}\\[4pt]
 &= \lim_{\Delta x\to 0}\frac{\textcolor{crimson}{\boldsymbol{b^{x}\cdot b^{\Delta x}}}-b^{x}}{\Delta x}
   &&\textcolor{crimson}{\text{\footnotesize\textbf{THE HINGE: available }\emph{only}\textbf{ because}}}\\[-3pt]
 & &&\textcolor{crimson}{\text{\footnotesize\textbf{the exponential maps sums}\,$\mapsto$\,\textbf{products}}}\\[4pt]
 &= \lim_{\Delta x\to 0}\frac{b^{x}\left(b^{\Delta x}-1\right)}{\Delta x} &&\text{\footnotesize group the common factor}\\[4pt]
 &= b^{x}\cdot\underbrace{\lim_{\Delta x\to 0}\frac{b^{\Delta x}-1}{\Delta x}}_{\textstyle \text{\normalsize a constant, call it }C}
   &&\text{\footnotesize $b^{x}$ contains no $\Delta x$, so it leaves the limit}\\[10pt]
 &= \;\boldsymbol{C\,b^{x}} &&
\end{align*}
\end{tcolorbox}
\end{center}

\begin{ideabox}[title={The intellectual centre of gravity of the whole Part}]
\begin{center}
\begin{tikzpicture}[font=\small]
 \node[rounded corners=4pt,fill=csB!15,draw=csB,line width=1pt,inner sep=8pt,text width=4.85cm,align=center]
   (A) at (-4.35,0) {maps \textbf{sums $\mapsto$ products}\\[2pt]{\scriptsize $b^{x+y}=b^{x}b^{y}$}};
 \node[rounded corners=4pt,fill=teal!12,draw=teal,line width=1pt,inner sep=8pt,text width=4.85cm,align=center]
   (B) at (4.35,0) {\textbf{rate of change $\propto$ size}\\[2pt]{\scriptsize $\frac{d}{dx}b^{x}=C\,b^{x}$}};
 \draw[{Stealth[length=7pt]}-{Stealth[length=7pt]},line width=1.4pt,draw=crimson] (A)--(B)
    node[midway,above=3pt,font=\scriptsize,text=crimson,align=center]{\textbf{mutually}\\\textbf{entailing}}
    node[midway,below=3pt,font=\scriptsize,text=crimson,align=center]{via one\\algebraic step};
\end{tikzpicture}
\end{center}
\vspace{-4pt}
A function that maps sums onto products \emph{necessarily} changes in proportion to its own size, and
conversely. The two ideas meet at a single line of algebra. Everything in Case Study 4 --- including
why $\cos y+i\sin y$ turns out to be an exponential at all --- runs through this junction.
\end{ideabox}

\begin{keybox}[title={The concept of $e$, stated before any number is computed}]
\begin{center}\large
\textit{``$e$'' is the base of the exponential function for which $C=1$:\\[4pt]
the function whose \textbf{rate of change is identical to itself}.}
\end{center}
\vspace{-2pt}
\hfill{\footnotesize\color{slate}\textsc{book p.~410--411 / pdf p.~429--430}}\\[4pt]
Equivalently, unpacked into the book's own summary box \pg{415}{434}: $e^{x}$ is the function that
\;\textbf{(i)} maps sums onto products, \;\textbf{(ii)} maps $1$ onto $2.718281828459045\ldots$,
\;\textbf{(iii)} is its own derivative, \;\textbf{(iv)} changes in exact proportion to its size.
\end{keybox}

\clearpage
\subsection{Why 2.718281828459045\ldots? The Power-Series blend}

Euler's own route to the number goes through power series, and the book's job is to make explicit the
blend that is taken for granted whenever a textbook differentiates a series term by term. It fuses
\textbf{six} domains.

\begin{center}\footnotesize
\setlength{\tabcolsep}{4pt}
\begin{tabular}{@{}>{\raggedright\arraybackslash}p{2.1cm}>{\raggedright\arraybackslash}p{2.1cm}>{\raggedright\arraybackslash}p{2.35cm}>{\raggedright\arraybackslash}p{2.6cm}>{\raggedright\arraybackslash}p{2.9cm}>{\raggedright\arraybackslash}p{3.1cm}@{}}
\toprule
\rowcolor{csB!16}
\textbf{\color{csB!75!black}1 Wholes \& parts} & \textbf{\color{csB!75!black}2 Numbers} &
\textbf{\color{csB!75!black}3 Functions} & \textbf{\color{csB!75!black}4 Infinite sums}\newline{\scriptsize via BMI} &
\textbf{\color{csB!75!black}5 Power series}\newline{\scriptsize via BMI} & \textbf{\color{csB!75!black}6 Calculus applied}\\
\midrule
Whole & resulting number $c$ & composite function $h(x)$ &
infinite sum $S=\sum_{k=0}^{\infty}a_k$ & function at infinity $F(x)=\sum a_k x^k$ &
derivative at infinity $F'(x)=\sum k\,a_k x^{k-1}$\\
\addlinespace[2pt]
Parts & operand numbers $a,b$ & component functions $f(x),g(x)$ &
finite terms $a_k$ & finite terms $a_k x^k$ & derivatives of finite terms $k\,a_k x^{k-1}$\\
\addlinespace[2pt]
part--whole relation & $a+b=c$ & $f(x)+g(x)=h(x)$ &
$a_{k-1}+a_k$ & $a_{k-1}x^{k-1}+a_kx^{k}$ & $(k{-}1)a_{k-1}x^{k-2}+k\,a_kx^{k-1}$\\
\addlinespace[2pt]
\rowcolor{csB!7}
--- & partial sum & partial sum of functions & $S_n=\sum_{k=0}^{n}a_k$ &
$F_n(x)=\sum_{k=0}^{n}a_kx^k$ & $F'_n(x)=\sum_{k=0}^{n}k\,a_kx^{k-1}$\\
\bottomrule
\end{tabular}
\end{center}
\vspace{-2pt}
\begin{center}\begin{minipage}{0.93\textwidth}\footnotesize\color{slate}
\textbf{Table 3.}\; The Power-Series blend \pg{413}{432}. Column 1$\to$2 is \meta{Arithmetic Is Object
Construction}; 2$\to$3 is \meta{Functions Are Numbers} (from Case Study 1); 2$\to$4 is the \textbf{BMI};
4$\to$5 is \meta{Functions Are Numbers} again, now applied to infinite sums; 6 applies differentiation
to an object that only exists because of columns 1--5.
\end{minipage}\end{center}

\subsubsection{The series is forced, not guessed}

Starting from ``$e^x$ is its own derivative'' plus $\frac{d}{dx}a_nx^{n}=n\,a_nx^{n-1}$, each term
must be the derivative of the term after it. And $x=0$ kills every term but the first, while
$e^{0}=1$, so $a_0=1$. Everything else cascades:

\begin{center}
\begin{tikzpicture}[font=\footnotesize,node distance=0pt,
  trm/.style={rounded corners=3pt,fill=csB!10,draw=csB,line width=0.8pt,inner sep=6pt,minimum height=0.95cm},
  bk/.style={-{Stealth[length=5pt]},line width=1pt,draw=crimson}]
\node[trm] (t0) at (0,0) {$1$};
\node[trm] (t1) at (2.1,0) {$x$};
\node[trm] (t2) at (4.5,0) {$\dfrac{x^{2}}{2!}$};
\node[trm] (t3) at (7.1,0) {$\dfrac{x^{3}}{3!}$};
\node[trm] (t4) at (9.7,0) {$\dfrac{x^{4}}{4!}$};
\node[font=\small] (dots) at (11.7,0) {$+\;\cdots$};
\foreach \a/\b in {t1/t0,t2/t1,t3/t2,t4/t3}{
  \draw[bk] (\a) to[out=140,in=40] node[midway,above=1pt,font=\scriptsize,text=crimson]{$\frac{d}{dx}$} (\b);}
\node[font=\scriptsize,text=slate,anchor=north,text width=13cm,align=center] at (5.8,-0.75)
  {each term differentiates \emph{backwards} onto the previous one --- so fixing the first term fixes
   the entire series};
\node[font=\scriptsize,text=csB!75!black,anchor=east,align=right,text width=2.25cm] at (-0.72,0) {$a_0=1$\\because $e^{0}=1$};
\end{tikzpicture}
\end{center}

\[
\boxed{\;e^{x}=1+x+\frac{x^{2}}{2!}+\frac{x^{3}}{3!}+\cdots\;}
\qquad\Longrightarrow\qquad
e=e^{1}=1+1+\frac{1}{2!}+\frac{1}{3!}+\cdots=2.718281828459045\ldots
\]

\begin{center}
\begin{tikzpicture}
\begin{axis}[width=8.6cm,height=5.4cm,
   xlabel={\footnotesize terms included, $n$}, ylabel={\footnotesize partial sum},
   xmin=0.5,xmax=8.5,ymin=1.8,ymax=2.95,
   xtick={1,2,3,4,5,6,7,8}, ytick={2.0,2.2,2.4,2.6,2.718281828,2.9},
   yticklabels={$2.0$,$2.2$,$2.4$,$2.6$,$e$,$2.9$},
   tick label style={font=\tiny}, label style={font=\footnotesize},
   grid=major, grid style={slate!18,line width=0.3pt},
   axis line style={slate!60},
   title={\footnotesize\bfseries\color{csB!75!black}$\sum_{k=0}^{n-1}1/k!\;\to\;e$}]
 \addplot[crimson,dashed,line width=1pt,domain=0.5:8.5,samples=2]{2.718281828};
 \addplot[csB,line width=1.2pt,mark=*,mark size=1.9pt,mark options={fill=csB}]
   coordinates {(1,1)(2,2)(3,2.5)(4,2.6666667)(5,2.7083333)(6,2.7166667)(7,2.7180556)(8,2.7182540)};
\end{axis}
\end{tikzpicture}
\hfill
\begin{tikzpicture}
\begin{axis}[width=8.6cm,height=5.4cm,
   xlabel={\footnotesize $n$}, ylabel={\footnotesize $(1+1/n)^{n}$},
   xmin=0,xmax=210,ymin=1.9,ymax=2.85,
   xtick={1,50,100,150,200}, ytick={2.0,2.2,2.4,2.6,2.718281828},
   yticklabels={$2.0$,$2.2$,$2.4$,$2.6$,$e$},
   tick label style={font=\tiny}, label style={font=\footnotesize},
   grid=major, grid style={slate!18,line width=0.3pt},
   axis line style={slate!60},
   title={\footnotesize\bfseries\color{csB!75!black}the textbook limit, converging \emph{much} more slowly}]
 \addplot[crimson,dashed,line width=1pt,domain=0:210,samples=2]{2.718281828};
 \addplot[teal,line width=1.2pt,domain=1:205,samples=120]{(1+1/x)^x};
\end{axis}
\end{tikzpicture}
\end{center}
\vspace{-4pt}
\begin{center}\begin{minipage}{0.93\textwidth}\footnotesize\color{slate}
\textbf{Figure 8.}\; Two routes to the same number \pg{414--419}{433--438}. The contrast is itself an
argument of the book's: the power series comes \emph{out of the concept} (``$e^x$ is its own
derivative''), whereas $(1+1/n)^{n}$ is normally handed to students as a definition with no
explanation. Note also how much faster the factorial series converges.
\end{minipage}\end{center}

\subsection{The appendix: the best-argued pages in Part VI}

\begin{warmbox}[title={The question textbooks never ask \normalfont\itshape(book p.~416 / pdf p.~435)}]
\begin{bookquote}{book p.~416 / pdf p.~435}
\textit{Why should this particular limit, $\lim_{n\to\infty}(1+\frac1n)^{n}$, be the base of the
exponential function that is its own derivative?}
\end{bookquote}
\end{warmbox}

\noindent Their answer is a five-link chain, and it starts from the \emph{meaning} of $e$, not from a
definition:

\begin{center}
\begin{tikzpicture}[font=\footnotesize,
  stp/.style={rounded corners=3pt,draw=csB,line width=0.9pt,fill=csB!8,inner sep=5pt,
              text width=2.72cm,align=center,minimum height=1.95cm},
  ar/.style={-{Stealth[length=5.5pt,width=4.5pt]},line width=1.1pt,draw=csB!75}]
\node[stp] (p1) at (0,0)     {$e^{x}$ is its own derivative, so the derivative at $x=1$ is $b$};
\node[stp] (p2) at (3.70,0)  {average change\\[3pt]$\dfrac{b^{1+q}-b}{q}\;\to\;b$};
\node[stp] (p3) at (7.40,0)  {use $b^{1+q}=b\cdot b^{q}$, factor out $b$:\\[3pt]$\dfrac{b^{q}-1}{q}=1$};
\node[stp] (p4) at (11.10,0) {solve for $b$:\\[3pt]$b^{q}=1+q$\\[3pt]$b=(1+q)^{1/q}$};
\node[stp,draw=crimson,fill=crimson!8] (p5) at (14.80,0)
   {set $q=\tfrac1n$:\\[3pt]$b=\lim\limits_{n\to\infty}\!\bigl(1+\tfrac1n\bigr)^{\!n}$};
\foreach \a/\b in {p1/p2,p2/p3,p3/p4,p4/p5}{\draw[ar] (\a)--(\b);}
\end{tikzpicture}
\end{center}
\vspace{-2pt}
\begin{center}\begin{minipage}{0.93\textwidth}\footnotesize\color{slate}
\textbf{Figure 9.}\; The derivation of the familiar limit \emph{from} the concept of $e$
\pg{416--419}{435--438}. The book's own summary of what the result rests on: ``(1) what $e$ means,
(2) the metaphors for mathematicizing change, (3) the limit metaphor, which is a special case of the
BMI, (4) some simple algebraic manipulations.''
\end{minipage}\end{center}

\begin{ideabox}[title={Their moral, and it is a good one}]
\begin{bookquote}{book p.~419 / pdf p.~438}
\textit{A system of conceptual metaphors and blends characterizes systematic relationships between
concepts and formulas. We have to understand the nature of that system if we are to understand the
mathematical ideas embodied by the formulas.}
\end{bookquote}
\end{ideabox}
